% =====================================================================
%  A dual linear programming bound for sphere packing in dimension 36
%  All numerical claims trace to code/receipt_d36_cutgen.txt,
%  code/receipt_d36_cs.txt and the exact ancillary data in data/.
% =====================================================================
\documentclass[11pt]{article}
\usepackage[margin=1.1in]{geometry}
\usepackage{amsmath,amssymb,amsthm}
\usepackage{booktabs}
\usepackage[colorlinks=true,linkcolor=blue,citecolor=blue,urlcolor=blue]{hyperref}

\emergencystretch=3em

\newtheorem{theorem}{Theorem}[section]
\newtheorem{lemma}[theorem]{Lemma}
\newtheorem{proposition}[theorem]{Proposition}
\newtheorem{corollary}[theorem]{Corollary}
\theoremstyle{remark}
\newtheorem{remark}[theorem]{Remark}

\newcommand{\LP}{\mathrm{LP}}
\newcommand{\Z}{\mathbb{Z}}
\newcommand{\R}{\mathbb{R}}
\newcommand{\HH}{\mathbb{H}}

\title{A dual linear programming bound for sphere packing\\ in dimension 36}
\author{Rifat Jumagulov\\ \texttt{jum.rifm@gmail.com}}
\date{July 13, 2026}

\begin{document}
\maketitle

\begin{abstract}
We construct an explicit dual-feasible point for the Cohn--Elkies linear
program in dimension $36$, built from the space of weight-$18$ modular
forms for $\Gamma_0(24)$ following the method of Cohn and Triantafillou.
The certificate shows that the two-point linear programming bound on the
sphere packing density in dimension $36$ exceeds the density of the best
packing currently known --- the Kschischang--Pasupathy packing, of center
density $2^{18}/3^{10}$ --- by a factor of at least $32.91$. In
particular, no Cohn--Elkies auxiliary function can certify the best known
packing in dimension $36$ as optimal. To our knowledge this is the first
such dual bound in any dimension above $32$, extending the table of
Cohn--Triantafillou ($d=12,16,20,28,32$), Li ($3\le d\le 13$), and
de~Courcy-Ireland--Dostert--Viazovska ($d=6$). The certificate is exact:
the dual point is a rational vector, coefficient nonnegativity is
verified by exact arithmetic up to $n=800$, and eventual positivity of
the two relevant $q$-expansions is proved via an explicit Deligne-type
tail bound whose constant is certified with outward-rounded interval
arithmetic. Two methodological points may be of independent interest: a
constraint-generation (cutting-plane) formulation of the exact rational
LP, which is what makes the tail-safe optimum reachable; and a sharpened,
lift-aware form of the Deligne bookkeeping constant,
$C=\sum_{f,e}|\lambda_{f,e}|\,e^{-(k-1)/2}$, without which the finite
verification in dimension $36$ fails marginally.
\end{abstract}

\section{Introduction}\label{sec:intro}

The Cohn--Elkies linear programming bound \cite{CohnElkies} is the
strongest known general upper bound on sphere packing densities in
moderate dimensions, and it is famously sharp in dimensions $8$ and $24$
\cite{Viazovska,CKMRV}. In every other dimension its sharpness is open,
and the question splits into two inequivalent problems. Write
$\delta_{\LP}(d)$ for the optimum of the Cohn--Elkies program
(normalized as a center density in the sense of \cite{ConwaySloane})
and $\delta_*(d)$ for the best packing density currently known.

\begin{itemize}
\item A \emph{dual bound} is a certified lower bound
  $\delta_{\LP}(d)\ge L$. If $L>\delta_*(d)$, then no Cohn--Elkies
  function can prove the optimality of the best known packing: either
  the LP bound is not sharp, or a denser packing than currently known
  exists. Cohn and Triantafillou \cite{CT} constructed such points from
  modular forms in dimensions $12,16,20,28,32$; Li \cite{Li} gave a
  different (discrete-reduction) construction covering $3\le d\le 13$;
  de~Courcy-Ireland, Dostert and Viazovska \cite{dCIDV} extended the
  modular-form method to dimension $6$.
\item \emph{Strict non-sharpness} is the stronger statement
  $\delta_{\LP}(d)>\delta_{\max}(d)$, where $\delta_{\max}$ is the true
  optimal density. This requires, in addition to a dual lower bound
  $L\le\delta_{\LP}(d)$, an \emph{independent} upper bound
  $U\ge\delta_{\max}(d)$ with $U<L$. At present this is known only for
  $d\in\{3,4,5,6,12,16\}$: for $d=12,16$ by combining the
  Cohn--Triantafillou dual points with the three-point bounds of
  Cohn--de~Laat--Salmon \cite{CdLS}, and for low dimensions by Li's dual
  bounds against classical upper bounds \cite{Li,dCIDV}.
\end{itemize}

This paper proves a statement of the first kind in dimension $36$, the
first above $32$.

\begin{theorem}\label{thm:main}
$\delta_{\LP}(36)\;\ge\;146.1036\ldots\;=\;32.9104\ldots\times
\dfrac{2^{18}}{3^{10}}.$
In particular the Cohn--Elkies linear programming bound in dimension
$36$ cannot certify the Kschischang--Pasupathy packing
\cite{KP}, of center density $2^{18}/3^{10}=4.4394\ldots$, as an optimal
sphere packing.
\end{theorem}

The witness is an explicit dual-feasible pair of modular forms
$(g,\tilde g)$ of weight $18$ for $\Gamma_0(24)$, with $\tilde g$ the
Fricke transform of $g$; the required nonnegativity of the two
$q$-expansions is verified exactly for $1\le n\le 800$ and proved for
all $n>n_0$ with an explicit crossover $n_0\le 63$, closing the gap.

Dimension $36$ is a natural target: the weight $k=d/2=18$ is even, so no
character twists are needed, and $36$ is the first dimension beyond the
Cohn--Triantafillou table where $\dim S_{d/2}$ obstructions do not
apply. It is also qualitatively different from the dimensions
$\le 32$: the ratio $\delta_{\LP}/\delta_*$ is much larger (our
certificate alone gives $32.9$; the two-point optimum is
$\approx 52\times\delta_*$), reflecting how weak the record packing is
relative to the LP bound in this range. Ironically this makes the
\emph{dual} statement stronger and the \emph{strict} question harder;
see \S\ref{sec:discussion}.

\paragraph{Why the construction is not routine.}
Running the Cohn--Triantafillou method at $d=36$ meets two obstacles
that do not occur (or are mild) at $d\le 32$.

First, the na\"ive LP relaxation (maximize the value subject to
nonnegativity of finitely many coefficients) produces optima whose
$q$-expansions turn negative just past the enforced range, on the
residue classes $n\equiv 3,9,15,21\ \  (\mathrm{mod}\ 24)$; while the
fully constrained problem (nonnegativity through the eventual-positivity
crossover, hundreds of rows with coefficients of size $\sim n^{17}$) is
numerically singular for floating-point solvers and prohibitively large
for exact dense simplex. We resolve this by \emph{constraint
generation}: solve the small relaxation exactly in rational arithmetic,
evaluate all violated rows exactly, add them, and re-solve. In dimension
$36$ this converges in a single round (all violations lie in the four
classes above and are eliminated together), producing an exact rational
vertex whose $q$-expansions are nonnegative far beyond the enforced
range (\S\ref{sec:construction}).

Second, the eventual-positivity verification of \cite[\S5]{CT} requires
an explicit constant $C$ with $|c_n|\le C\,\sigma_0(n)\,n^{(k-1)/2}$ for
the cuspidal parts of \emph{both} $g$ and $\tilde g$. The natural
constant --- the $\ell^1$-norm $\sum|\lambda_{f,e}|$ of the coordinates
in the newform-lift basis --- turns out to be \emph{lossy} in a way that
matters here: our certificate's cuspidal part is dominated ($99.7\%$ of
the norm) by a single deep oldform lift, the $e=24$ shift of the level-one
newform, and the na\"ive constant fails the required threshold
marginally. The lift-aware constant
$C_w=\sum_{f,e}|\lambda_{f,e}|\,e^{-(k-1)/2}$
(Lemma~\ref{lem:liftaware}), equally rigorous and elementary, is smaller
here by ten orders of magnitude and closes the verification with a
margin of $7\times 10^8$ (\S\ref{sec:tail}).

A further point, absent from dimensions where the certificate has a
sparse structure, is that \emph{both} cusps genuinely matter: not only
the Fricke transform $\tilde g$ but also $g$ itself has a razor-thin
Eisenstein component ($e_1(g)\approx -3.9\times10^{-11}$), so the
$a_n\ge 0$ side needs the same Deligne-type argument as the
$b_n\ge 0$ side, with its own constants (\S\ref{sec:tail}).

\paragraph{Rigor.} The dual vertex is an exact rational vector (the
entries have $\approx 128$-digit numerators). All coefficient scans up
to $n=800$ (and, in a second implementation, to $n=1100$) are exact.
The newform decompositions of both cuspidal parts are validated by an
exact identity: the reconstruction of the $q$-expansion from the
computed coordinates agrees with the exact coefficients \emph{as
rational numbers} (deviation identically $0$) up to $n=800$. The tail
constants $C_w$ are certified as upper bounds using outward-rounded
interval arithmetic, with the algebraic Hecke eigenvalues enclosed by
exact Sturm isolating intervals. Independent implementations (exact
characteristic-polynomial factorization vs.\ high-precision simultaneous
diagonalization) produce identical constants. All code, receipts, and
the exact certificate data are provided as ancillary files (see the
manifest in \texttt{ancillary/MANIFEST.md}).

\section{The Cohn--Elkies program and its modular-form dual}
\label{sec:dual}

The engine is the following observation of Cohn and Triantafillou
\cite[Prop.~2.1]{CT}: if $\mu$ is a tempered distribution on $\R^d$
with $\mu=\delta_0+\nu$, $\nu\ge0$,
$\operatorname{supp}(\nu)\subseteq\{x\in\R^d:|x|\ge r\}$ for some
$r>0$, and $\hat\mu\ge c\,\delta_0$ for some $c>0$, then the
Cohn--Elkies linear programming bound in $\R^d$ is at least
$c\,(r/2)^d$. (Any Cohn--Elkies auxiliary function $f$ pairs with $\mu$
on both sides of the inequality chain, forcing
$f(0)/\hat f(0)\ge c\,r^d$ up to normalization.) Sections 2--3 of
\cite{CT} construct such distributions from modular forms, as follows.

Fix the dimension $d=36$, weight $k=d/2=18$, and level $N=24$. For a
holomorphic modular form $g\in M_k(\Gamma_0(N))$ write
$g=\sum_{n\ge0}a_n q^n$ ($q=e^{2\pi iz}$) and let
\[
\tilde g \;=\; i^{k}\,\bigl(g\big|_k w_N\bigr),
\qquad
w_N=\begin{pmatrix}0&-1\\ N&0\end{pmatrix},
\qquad
(f|_k M)(z)=(\det M)^{k/2}(cz+d)^{-k}f(Mz),
\]
be its (normalized) Fricke transform, i.e.\
$\tilde g(z)=i^kN^{-k/2}z^{-k}g(-1/(Nz))$; for our even weight $k=18$
the factor is $i^{18}=-1$. Then
$\tilde g\in M_k(\Gamma_0(N))$ as well, and
$\tilde{\tilde g}=g$. Write $\tilde g=\sum_{n\ge0}b_nq^n$. Given $g$
satisfying the coefficient conditions below, the construction of
\cite[\S\S2--3]{CT} produces a distribution $\mu$ as above with
$r=\sqrt T$ and $c=(2/\sqrt N)^{d/2}\,b_0$. Concretely: from any $g$
satisfying
\begin{equation}\label{eq:feas}
a_0=1,\qquad a_n=0\ (1\le n\le T-1),\qquad a_n\ge 0\ (n\ge T),
\qquad b_n\ge 0\ (n\ge 0),
\end{equation}
a dual-feasible point for the Cohn--Elkies program with value (as a
center density)
\begin{equation}\label{eq:value}
B(g)\;=\;b_0\cdot\Bigl(\tfrac{2}{\sqrt N}\Bigr)^{d/2}\cdot
\Bigl(\tfrac{\sqrt T}{2}\Bigr)^{d},
\end{equation}
so that $\delta_{\LP}(d)\ge B(g)$. Here $T$ is a free integer parameter
(the number of forced initial zeros). We take $T=10$. As a consistency
check of the normalization, our implementation of
\eqref{eq:feas}--\eqref{eq:value} reproduces the published $d=12$
values of \cite{CT} at both levels: $0.059782=1.6141\times$
Coxeter--Todd at $N=24$ (\cite[Table~6.1]{CT}) and
$0.0624462=1.68605\times$ at $N=96$ (\cite[Thm.~1.1]{CT}), to the
displayed precision.

\section{The certificate}\label{sec:construction}

\subsection{The space and the reduced LP}

$\dim M_{18}(\Gamma_0(24))=72$, with an $8$-dimensional Eisenstein
subspace spanned by $E_{18}(\delta z)$, $\delta\mid 24$, where
\[
E_{18}(z)=1+\kappa_{18}\sum_{n\ge1}\sigma_{17}(n)q^n,
\qquad
\kappa_{18}=-\frac{2k}{B_{18}}\Big|_{k=18}=-\frac{28728}{43867},
\]
and a $64$-dimensional cuspidal subspace. A spanning set of $72$
linearly independent forms is obtained from the eight Eisenstein series
together with holomorphic eta-quotients of weight $18$ and level
dividing $24$ (Ligozat's criteria); all $q$-expansions are computed in
exact integer arithmetic. Imposing the $9$ linear conditions
$a_1=\dots=a_9=0$ and the normalization $a_0=1$ leaves an affine space
of dimension $r=20$, on which we solve
\begin{equation}\label{eq:lp}
\max\ b_0
\quad\text{s.t.}\quad
a_n\ge0\ (10\le n\le 28),\qquad
b_n\ge0\ (n\in W),
\end{equation}
with a working set $W$ of Fricke-side rows, in exact rational
arithmetic.

\subsection{Constraint generation}

With $W=\{1,\dots,28\}$ the exact optimum of \eqref{eq:lp} is
$36.3470\times 2^{18}/3^{10}$, but its Fricke expansion has $45$
negative coefficients in $1\le n\le 300$, all in the classes
$n\equiv3,9,15,21\  (\mathrm{mod}\ 24)$, the first at $n=33$: the
unconstrained optimum overshoots into infeasibility, and (as we show
below) its Eisenstein data forces $b_n<0$ for infinitely many $n$, so no
finite verification can rescue it. Adding all $45$ violated rows to $W$
($|W|=73$) and re-solving exactly gives a vertex with
\[
b_0 \;=\; 101.181761\ldots \text{ (exact rational)},\qquad
B \;=\; 146.10367\ldots \;=\; 32.9104\ldots\times\frac{2^{18}}{3^{10}},
\]
and \emph{no} violations: re-scanning exactly, $a_n\ge0$ and $b_n\ge0$
for all $1\le n\le 800$ (zero negative coefficients on either side,
including the un-enforced range $300<n\le800$). The cutting-plane
process converges in one round. Exact rational arithmetic in the solve
is essential: at this conditioning, floating-point and even
$300$-digit fixed-precision simplex return phantom vertices (violating
$a_0=1$) once rows with coefficients of size $10^{30}$ enter the basis.

\begin{remark}
The certificate data --- the $20$ exact rational coordinates of the
vertex (in the reduced basis of the construction), the exact Eisenstein
coefficients, and all constants below --- are in the ancillary file
\texttt{ancillary/certificate\_exact\_data.txt}; the construction and
verification receipts are \texttt{ancillary/receipt\_d36\_cutgen.txt}
and \texttt{ancillary/receipt\_d36\_cs.txt}.
\end{remark}

\section{Eventual positivity}\label{sec:tail}

It remains to prove $a_n\ge0$ and $b_n\ge0$ for all $n>800$. Decompose
both forms into Eisenstein and cuspidal parts; we treat $\tilde g$ (the
binding side) and note the same argument for $g$ with its own constants.

\subsection{Eisenstein main term, per residue class}

Write the Eisenstein component of $\tilde g$ as
$\sum_{\delta\mid24}e_\delta\,E_{18}(\delta z)$. The exact rational
$e_\delta$ (computed by an exact projector, validated as explained
below) have the numerical values
\[
\begin{array}{ll}
e_1=-4.197922\cdot10^{-13}, & e_2=-2.023981\cdot10^{-14}, \\
e_3=+5.361611\cdot10^{-5}, & e_4=+5.776815\cdot10^{-3}, \\
e_6=-18.24676, & e_8=-27.55171, \\
e_{12}=+43.87748, & e_{24}=+103.0969.
\end{array}
\]
For $n\ge1$ the Eisenstein coefficient is
$\kappa_{18}\sum_{\delta\mid\gamma}e_\delta\,\sigma_{17}(n/\delta)$,
where $\gamma=\gcd(n,24)$. Using
$m^{17}\le\sigma_{17}(m)<\zeta(17)\,m^{17}$ together with the rational
bound $\zeta(17)<1+2^{-17}+2^{-16}/16$ (compare the sum with
$\int_2^\infty x^{-17}dx$), and bounding each term from below according
to its sign (note $\kappa_{18}<0$, so $\kappa_{18}e_\delta>0$ exactly
when $e_\delta<0$), we get for every $n\ge1$ in the class $\gamma$:
\begin{equation}\label{eq:guard}
\text{Eisenstein}(n)\;\ge\;c_E(\gamma)\,n^{17},
\qquad
c_E(\gamma)\;=\;
\sum_{\substack{\delta\mid\gamma\\ e_\delta<0}}
\frac{\kappa_{18}e_\delta}{\delta^{17}}
\;+\;
\Bigl(1+2^{-17}+\tfrac{2^{-16}}{16}\Bigr)
\sum_{\substack{\delta\mid\gamma\\ e_\delta>0}}
\frac{\kappa_{18}e_\delta}{\delta^{17}},
\end{equation}
an exact rational number for each $\gamma$. Evaluating exactly: all
eight constants $c_E(\gamma)$, $\gamma\mid24$, are \emph{positive},
with the minimum on the class $\gamma=3$:
\begin{equation}\label{eq:cE}
c_E \;:=\; \min_{\gamma\mid24}c_E(\gamma)\;=\;c_E(3)
\;=\;3.0197203\cdot10^{-15}\quad(\tilde g\text{-side}),
\end{equation}
and analogously $c_E^{(a)}=c_E^{(a)}(3)=5.0672218\cdot10^{-12}$ for
$g$ (again the minimum over all eight classes, all positive). The
displayed decimals are certified truncations of the exact rationals,
which are provided in the ancillary files.
The thinness of these constants --- forced by the active constraints at
the LP optimum --- is what makes the tail argument in dimension $36$
delicate: the certificate value lives $12$--$14$ orders of magnitude
above the Eisenstein margin that protects its tail.

\subsection{The cuspidal remainder and a lift-aware Deligne constant}

Let $S=\tilde g-\text{(Eisenstein part)}=\sum c_nq^n\in
S_{18}(\Gamma_0(24))$, a space of dimension $64$. The newform/oldform
basis consists of the lifts $f(ez)$, where $f$ runs over the normalized
newforms of level $M\mid24$ and $e\mid 24/M$; the newspace dimensions
are
\[
(\dim S^{\mathrm{new}}_{18}(\Gamma_0(M)))_{M\mid24}
=(1,1,3,2,3,4,2,9)\ \text{ for }M=(1,2,3,4,6,8,12,24),
\]
with lift multiplicities $(8,6,4,4,3,2,2,1)$, summing to $64$. We
compute the exact coordinates $\lambda_{f,e}$ of $S$ in this basis
(Galois orbits handled by exact characteristic-polynomial factorization
of the Hecke action; bad-prime lifts separated by $U_2,U_3$ and leading
$q$-exponents) and \emph{validate the decomposition exactly}: the
reconstruction $\sum\lambda_{f,e}f(ez)$ agrees with $S$
coefficient-by-coefficient, as exact rationals, for all $n\le 800$
(deviation identically zero; the Sturm bound for
$M_{18}(\Gamma_0(24))$ is $72$, so agreement to $800$ is a highly
overdetermined identity). As a further check, the level-one component is
the unique normalized cusp form $\Delta E_6$ of weight 18, and the computed
eigenvalues match $a_2=-528$, $a_3=-4284$, $a_5=-1025850$,
$a_7=3225992$ exactly.

\begin{lemma}[lift-aware Deligne constant]\label{lem:liftaware}
Let $S=\sum_{f,e}\lambda_{f,e}\,f(ez)\in S_k(\Gamma_0(N))$ with $f$
normalized newforms. Then for all $n\ge1$,
\[
|c_n|\;\le\;C_w\,\sigma_0(n)\,n^{(k-1)/2},
\qquad
C_w\;=\;\sum_{f,e}\frac{|\lambda_{f,e}|}{e^{(k-1)/2}}.
\]
\end{lemma}

\begin{proof}
By Deligne's bound $|a_f(m)|\le\sigma_0(m)m^{(k-1)/2}$ for each
normalized newform. The coefficient of $q^n$ in $f(ez)$ is $a_f(n/e)$
if $e\mid n$ and $0$ otherwise, and for $e\mid n$,
$\sigma_0(n/e)\le\sigma_0(n)$ and
$(n/e)^{(k-1)/2}=n^{(k-1)/2}e^{-(k-1)/2}$. Summing over the pairs
$(f,e)$ with $e\mid n$ gives the claim.
\end{proof}

The point of the lemma is quantitative. The na\"ive constant
$\sum_{f,e}|\lambda_{f,e}|$ equals $1.894\times10^{10}$ for our $S$:
it is dominated ($99.7\%$) by the single coordinate
$\lambda_{\Delta E_6,\,24}\approx-1.89\times10^{10}$ on the deepest
oldform lift, and with it the crossover lands just \emph{beyond} the
verified range ($n_0\approx827>800$) --- the verification fails
marginally. The lift-aware constant suppresses that term by
$24^{17/2}\approx5.4\times10^{11}$:
\begin{equation}\label{eq:Cw}
C_w \;\le\; 0.358807778339 \quad(\tilde g),
\qquad
C_w^{(a)} \;\le\; 0.306329333530 \quad(g),
\end{equation}
certified as upper bounds by outward-rounded interval arithmetic: the
coordinates on rational (one-dimensional Galois orbit) blocks, including
the dominant level-one tower, are computed as exact rational numbers
(31 of the 64 lifts); the remaining 33 algebraic coordinates are
enclosed via exact Sturm isolating intervals for the Hecke eigenvalues
and an interval linear solve (enclosure widths $\sim10^{-50}$; the bound
is stable to $15$ digits across four disjoint determining windows, and
degrading the working precision widens the enclosure without breaking
the bound, confirming outward rounding).

\subsection{Closing the certificate}

Combining \eqref{eq:cE}, Lemma~\ref{lem:liftaware}, \eqref{eq:Cw} and
$\sigma_0(n)\le2\sqrt n$: for $n$ on any residue class,
\[
b_n\;\ge\;c_E\,n^{17}-2\,C_w\,n^{9}\;>\;0
\qquad\text{whenever } n^{8}>\frac{2C_w}{c_E},
\]
i.e.\ for all $n\ge n_0=\lceil(2C_w/c_E)^{1/8}\rceil=63$; and on the
$g$-side for all $n\ge n_0^{(a)}=25$. Since $a_n\ge0$ and $b_n\ge0$ have
been verified exactly for $1\le n\le800$, both hold for all $n$, and
$(g,\tilde g)$ satisfies \eqref{eq:feas}. Theorem~\ref{thm:main}
follows from \eqref{eq:value}. \qed

The margins are enormous ($2C_w/c_E$ would have to exceed its actual
value by a factor $7\times10^{8}$ before $n_0$ reached $800$); the
certificate is in no sense borderline once the correct constant is used.

\section{Discussion}\label{sec:discussion}

\paragraph{Strict non-sharpness in dimension 36 is far out of reach.}
Our dual point shows $\delta_{\LP}(36)\ge32.91\,\delta_*$, and the
two-point optimum itself is $\approx52.2\,\delta_*$
\cite{CohnElkies,CT}. Strict non-sharpness would need an independent
upper bound below $32.91\,\delta_*$, i.e.\ an improvement of
$\approx37\%$ over the two-point bound. For comparison, the only
technology that has ever improved on the two-point bound --- the
three-point bounds of \cite{CdLS} --- gains $4.85\%$ at $d=12$ and
$0.79\%$ at $d=16$, and no three-point computation exists in any
dimension beyond $16$ \cite{CdLS,VallentinSurvey}. The paradox of
dimension $36$ is that the \emph{weakness} of the known record makes
the dual statement strong and the strict statement inaccessible: the
dual value sits far below the LP optimum precisely because the record
is far below both.

\paragraph{Dimensions 40 and 44.} The same pipeline reaches $d=40$
(weight $20$) in principle, with two new costs: the weight-$20$
eta-quotients of level $24$ span only a $72$-dimensional subspace of the
$80$-dimensional $M_{20}(\Gamma_0(24))$, so the basis must be augmented
(e.g.\ by $E_4\cdot(\text{weight-16 forms})$), and the exact arithmetic
is heavier ($\sigma_{19}$ growth). We have not completed a $d=40$
certificate. We note that Remark~4.3 of \cite{Zhou} asserts that for
$d=32$ and $d=40$ ``LP non-sharpness is established by other means'',
without a supporting citation for the case $d=40$; we are not aware of
any dual bound or non-sharpness proof in dimension $40$. The
modular-bootstrap analysis of \cite{AJCHdLT} (building on the sphere
packing--quantum gravity correspondence of \cite{HMR}) rules out
\emph{sharpness}
of the LP bound for a range of dimensions by an implied-kissing-number
criterion, but reports its first crossings of that criterion only near
$d\approx180$ and, in its authors' own words, ``does not give a
quantitative improvement''; it therefore neither contains nor implies a
dual bound at $d=36$ or $40$.

\paragraph{Methodological remarks.} Two ingredients may be reusable.
(i) Constraint generation appears to be the right formulation of the
exact dual LP whenever eventual positivity must be enforced through a
crossover: it keeps every exact solve small (here: one re-solve with
$73$ rows) while the dense exact LP is intractable and floating-point
LP is unsound at these conditionings. (ii) Lemma~\ref{lem:liftaware} is
elementary, but the failure of the unweighted constant in dimension
$36$ suggests it should be the default bookkeeping in
eventual-positivity verifications: deep oldform lifts are exactly the
coordinates the unweighted $\ell^1$ norm overweights, by the factor
$e^{(k-1)/2}$ that the lemma restores.

\paragraph{Reproducibility.} The exact verification --- the rigorous part
of the certificate --- uses only exact rational linear algebra, integer
$q$-expansions of eta-quotients, and interval arithmetic (Python:
\texttt{sympy}, \texttt{mpmath}, \texttt{fractions}); no computer-algebra
system with modular-forms machinery is required. Floating-point routines
(\texttt{numpy}, \texttt{scipy}) are used only in the \emph{construction}
step, as a heuristic to locate the candidate vertex and basis; they play
no role in the exact rigor pass, and every constant reported here is
recomputed exactly. The full pipeline, receipts, and the exact
certificate are provided as ancillary files, together with an
independent second implementation of the newform decomposition
(high-precision simultaneous diagonalization) that reproduces the
constants \eqref{eq:Cw} to all reported digits.

\paragraph{Data and code availability.} The exact certificate data, the
construction and verification receipts, and the complete Python pipeline
(including the independent second implementation) are provided as
ancillary files accompanying this submission; see
\texttt{ancillary/MANIFEST.md} for the file map and reproduction order.

\paragraph*{Acknowledgements.} The construction, the exact-arithmetic
certificate, and the preparation of this manuscript were carried out with
substantial assistance from AI language-model tools. The
correctness of the main result does not rest on trust in those tools: the
certificate is exact and machine-checkable, and every numerical claim has
been reproduced by re-running the exact-arithmetic verification from
scratch (the pipeline and receipts are provided as ancillary files). The
author is accountable for this work and for the decision to publish it.

\subsection*{Mathematics Subject Classification (2020)}
Primary 52C17 (Packing and covering in $n$ dimensions);
Secondary 11F11 (Holomorphic modular forms of integral weight),
11F30 (Fourier coefficients of automorphic forms),
90C05 (Linear programming).

\medskip\noindent\emph{Keywords:} sphere packing, linear programming
bound, Cohn--Elkies bound, dual bound, modular forms, eta-quotients,
eventual positivity, exact arithmetic.

\begin{thebibliography}{99}

\bibitem{CohnElkies} H.~Cohn and N.~Elkies,
\emph{New upper bounds on sphere packings I},
Ann. of Math. \textbf{157} (2003), 689--714. arXiv:math/0110009.

\bibitem{Viazovska} M.~Viazovska,
\emph{The sphere packing problem in dimension 8},
Ann. of Math. \textbf{185} (2017), 991--1015.

\bibitem{CKMRV} H.~Cohn, A.~Kumar, S.~D.~Miller, D.~Radchenko and
M.~Viazovska, \emph{The sphere packing problem in dimension 24},
Ann. of Math. \textbf{185} (2017), 1017--1033.

\bibitem{CT} H.~Cohn and N.~Triantafillou,
\emph{Dual bounds for sphere packing via modular forms},
Math. Comp. \textbf{91} (2022), 491--508. arXiv:1909.04772.

\bibitem{Li} R.~Li,
\emph{Dual linear programming bounds for sphere packing via discrete
reductions}, Adv. Math. \textbf{460} (2024), 110043. arXiv:2206.09876.

\bibitem{dCIDV} M.~de~Courcy-Ireland, M.~Dostert and M.~Viazovska,
\emph{Six-dimensional sphere packing and linear programming},
Math. Comp. \textbf{93} (2024), 1993--2029. arXiv:2211.09044.

\bibitem{CdLS} H.~Cohn, D.~de~Laat and A.~Salmon,
\emph{Three-point bounds for sphere packing}. arXiv:2206.15373.

\bibitem{KP} F.~R.~Kschischang and S.~Pasupathy,
\emph{Some ternary and quaternary codes and associated sphere packings},
IEEE Trans. Inform. Theory \textbf{38} (1992), 227--246.

\bibitem{AJCHdLT} N.~Afkhami-Jeddi, H.~Cohn, T.~Hartman, D.~de~Laat and
A.~Tajdini, \emph{High-dimensional sphere packing and the modular
bootstrap}, J. High Energy Phys. (2020). arXiv:2006.02560.

\bibitem{HMR} T.~Hartman, D.~Maz\'a\v{c} and L.~Rastelli,
\emph{Sphere packing and quantum gravity},
J. High Energy Phys. (2019). arXiv:1905.01319.

\bibitem{Zhou} J.~Zhou,
\emph{Cusp form dimensions, lattice uniqueness, and LP sharpness for
sphere packing in dimensions 8 and 24}. arXiv:2604.10914.

\bibitem{VallentinSurvey} F.~Vallentin,
\emph{Conic optimization for extremal geometry}. arXiv:2510.06960.

\bibitem{ConwaySloane} J.~H.~Conway and N.~J.~A.~Sloane,
\emph{Sphere Packings, Lattices and Groups}, 3rd ed., Springer, 1999.

\end{thebibliography}

\end{document}
